\subsubsection{异常类的一点判别与交叉异常坐标化：torsion-free 判据与 Euler 商不变量}\label{subsubsec:pom-anom-torsionfree-euler-quotient}
本段在三生成元片段的异常上同调口径（推论 \ref{cor:pom-anom-cohomology-class}）与其 moment 放大律（命题 \ref{prop:pom-anom-moment-amplification}）之上，
给出一组与常数层审计直接相关的判别与坐标化结论：
其一在 torsion-free 情形下给出\emph{一点零异常判别}；
其二在不预设 torsion-free 的情形下给出\emph{互素两点完备判别}，并由此推出“全矩阶零异常”的两点有限化（预言机坍缩）；
其三刻画一般交换群上的\emph{torsion 可消去性与最小消去阶}，并在环面型目标群上给出分母的最小公倍数律；
其四把高阶交叉异常族 \(\{\Delta_S\}\) 组织为 Euler/直积可分解性的\emph{最小线性坐标}（商空间同构与维数公式）。

\begin{corollary}[torsion-free 情形的一点零异常判别]\label{cor:pom-anom-one-point-zero-test}
令 \(H\) 表示投影词片段的异常上同调群（由中心标签群按 coboundary 取商得到的交换群），并令 \([\mathrm{Anom}(w)]\in H\) 为任意 \(1\)-态射 \(w\) 的异常类。
若 \(H\) 为 torsion-free（即 \(qa=0\Rightarrow a=0\) 对一切整数 \(q\ge 1\) 成立），则对任意整数 \(q\ge 2\) 有严格等价
\[
\boxed{\;
[\mathrm{Anom}(w)]=0\quad\Longleftrightarrow\quad [\mathrm{Anom}(\mathrm{MOM}_q\circ w)]=0\; }.
\]
因此在该类 \(H\) 上，”零异常”可由任意单个矩阶的核验完全决定。
\leanverified{torsion\_free\_smul\_eq\_zero\_iff}
\end{corollary}
\begin{proof}
由命题 \ref{prop:pom-anom-moment-amplification}，
\(
[\mathrm{Anom}(\mathrm{MOM}_q\circ w)]=q\,[\mathrm{Anom}(w)]
\)。
若 \([\mathrm{Anom}(w)]=0\) 则右端为零，故推出 \([\mathrm{Anom}(\mathrm{MOM}_q\circ w)]=0\)。
反之若 \([\mathrm{Anom}(\mathrm{MOM}_q\circ w)]=0\)，则 \(q\,[\mathrm{Anom}(w)]=0\)，由 torsion-free 得 \([\mathrm{Anom}(w)]=0\)。证毕。
\end{proof}

\begin{theorem}[torsion-free 情形下算法等价的单矩完备判别]\label{thm:pom-anom-single-moment-complete-equivalence}
\leanverified{paper\_pom\_anom\_single\_moment\_complete\_equivalence}
在命题 \ref{prop:pom-confluence-mod-anom-audited} 与定义 \ref{def:pom-algorithmic-equivalence-nf-anom} 的有界审计语义下，
设异常上同调群 \(H\) 为 torsion-free，并固定任意整数 \(q\ge 2\)。
则对任意两个投影词 \(u,v\)，有严格等价
\[
\boxed{\;
u\sim v
\quad\Longleftrightarrow\quad
\mathrm{NF}(u)=\mathrm{NF}(v)
\ \text{且}\
[\mathrm{Anom}(\mathrm{MOM}_q\circ u)]
=
[\mathrm{Anom}(\mathrm{MOM}_q\circ v)]\; }.
\]
换言之，在正规形一致的前提下，算法等价的异常核验可压缩为任意单个非平凡矩阶上的异常切片比较。
\end{theorem}
\begin{proof}
由定义 \ref{def:pom-algorithmic-equivalence-nf-anom}，
\[
u\sim v
\quad\Longleftrightarrow\quad
\mathrm{NF}(u)=\mathrm{NF}(v)
\ \text{且}\
[\mathrm{Anom}(u)]=[\mathrm{Anom}(v)].
\]
在 \(\mathrm{NF}(u)=\mathrm{NF}(v)\) 的前提下，令
\[
a:=[\mathrm{Anom}(u)]-[\mathrm{Anom}(v)]\in H.
\]
则 \(u\sim v\) 等价于 \(a=0\)。
另一方面，由命题 \ref{prop:pom-anom-moment-amplification}，
\[
[\mathrm{Anom}(\mathrm{MOM}_q\circ u)]
-
[\mathrm{Anom}(\mathrm{MOM}_q\circ v)]
=
q\,a.
\]
因此“\(q\)-阶 moment 切片上的异常类相等”等价于 \(q\,a=0\)。
由于 \(H\) 为 torsion-free，\(q\,a=0\) 当且仅当 \(a=0\)；与上式合并即得结论。证毕。
\end{proof}

\begin{theorem}[互素矩阶零消失原理：异常检验的两点完备性（无需 torsion-free）]\label{thm:pom-anom-coprime-two-point-test}
设 PW 语义下的异常类落在某交换群 \(H\) 中，并满足矩放大律：对任意整数 \(q\ge 2\)，有
\[
[\mathrm{Anom}(\mathrm{MOM}_q\circ w)]=q\cdot[\mathrm{Anom}(w)]\in H
\]
（命题 \ref{prop:pom-anom-moment-amplification}）。
若存在互素整数 \(q_1,q_2\ge 2\) 使得
\[
[\mathrm{Anom}(\mathrm{MOM}_{q_1}\circ w)]=0,\qquad
[\mathrm{Anom}(\mathrm{MOM}_{q_2}\circ w)]=0,
\]
则必有 \([\mathrm{Anom}(w)]=0\)。
从而对所有 \(q\ge 2\) 均成立 \([\mathrm{Anom}(\mathrm{MOM}_{q}\circ w)]=0\)。
\leanverified{coprime\_smul\_eq\_zero\_of\_both}
\end{theorem}
\begin{proof}
令 \(a:=[\mathrm{Anom}(w)]\in H\)。由放大律得 \(q_1 a=0\) 与 \(q_2 a=0\)。
取整数 \(u,v\) 使得 \(u q_1+v q_2=1\)（Bézout 恒等式），则
\[
a=(u q_1+v q_2)a=u(q_1 a)+v(q_2 a)=0.
\]
再代回放大律即得对一切 \(q\ge 2\) 的结论。证毕。
\end{proof}

\begin{theorem}[异常—预言机坍缩：无限矩阶审计 \texorpdfstring{$\Longleftrightarrow$}{<->} 两个互素矩阶审计]\label{thm:pom-anom-all-moments-collapse-to-two}
令命题
\[
\mathcal P(w):\quad \forall q\ge 2,\;[\mathrm{Anom}(\mathrm{MOM}_q\circ w)]=0
\]
表示“异常在全矩阶上消失”。
则 \(\mathcal P(w)\) 与存在互素整数 \(q_1,q_2\ge 2\) 使
\(
[\mathrm{Anom}(\mathrm{MOM}_{q_1}\circ w)]=[\mathrm{Anom}(\mathrm{MOM}_{q_2}\circ w)]=0
\)
等价。
\leanverified{anom\_oracle\_collapse}
\end{theorem}
\begin{proof}
“\(\Rightarrow\)”显然（例如取 \((q_1,q_2)=(2,3)\)）。
“\(\Leftarrow\)”由定理 \ref{thm:pom-anom-coprime-two-point-test} 立得。证毕。
\end{proof}

\begin{corollary}[预言机坍缩的可审计触发条件：全矩阶零异常的两点有限证书]\label{cor:pom-anom-oracle-collapse-two-point-trigger}
\leanverified{paper\_pom\_anom\_oracle\_collapse\_two\_point\_trigger}
凡以命题 \(\mathcal P(w)\) 作为充分条件、从而在逻辑上依赖“对所有矩阶 \(q\ge 2\)”的无限审计者，
其触发条件均可压缩为仅核验两个互素矩阶；换言之，“全矩阶异常轴”的审计复杂度在逻辑上坍缩为两点。
\end{corollary}

\begin{theorem}[torsion 可消去性判别与最小消去阶]\label{thm:pom-anom-torsion-eliminability-min-order}
在定理 \ref{thm:pom-anom-coprime-two-point-test} 的设定下，令 \(a:=[\mathrm{Anom}(w)]\in H\)，并记
\[
H_{\mathrm{tor}}:=\{x\in H:\ \exists\,n\ge 1\ \text{使}\ nx=0\}
\]
为 \(H\) 的 torsion 子群。
则：
\begin{enumerate}
  \item 存在整数 \(q\ge 2\) 使 \([\mathrm{Anom}(\mathrm{MOM}_q\circ w)]=0\) 当且仅当 \(a\in H_{\mathrm{tor}}\)。
  \item 若 \(a\notin H_{\mathrm{tor}}\) 且 \(H\) 可嵌入为某赋范实线性空间 \((V,\|\cdot\|)\) 的加法子群，则
  \[
  \bigl\|[\mathrm{Anom}(\mathrm{MOM}_q\circ w)]\bigr\|
  =\|qa\|\xrightarrow[q\to\infty]{}\infty.
  \]
  \item 若 \(a\in H_{\mathrm{tor}}\) 且其阶为 \(\mathrm{ord}(a)\in\NN\)，则对任意整数 \(q\ge 2\) 有严格等价
  \[
  \boxed{\;
  [\mathrm{Anom}(\mathrm{MOM}_q\circ w)]=0\quad\Longleftrightarrow\quad \mathrm{ord}(a)\mid q\; }.
  \]
  特别地，\(\mathrm{ord}(a)\) 为异常可由矩编译消去的最小正整数（最小消去阶）。
\end{enumerate}
\leanverified{smul\_eq\_zero\_iff\_order\_dvd}
\end{theorem}
\begin{proof}
由命题 \ref{prop:pom-anom-moment-amplification}，
\(
[\mathrm{Anom}(\mathrm{MOM}_q\circ w)]=qa
\)。
故存在 \(q\ge 2\) 使其为零当且仅当存在 \(q\ge 2\) 使 \(qa=0\)，这等价于 \(a\in H_{\mathrm{tor}}\)。
若 \(a\notin H_{\mathrm{tor}}\) 且 \(H\hookrightarrow (V,\|\cdot\|)\)，则 \(a\neq 0\) 且 \(\|qa\|=q\|a\|\to\infty\)。
若 \(a\in H_{\mathrm{tor}}\) 且 \(\mathrm{ord}(a)=n\)，则 \(qa=0\) 当且仅当 \(n\mid q\)，从而最小正整数即为 \(n\)。证毕。
\end{proof}

\begin{corollary}[环面型相位异常的有限覆盖消去与分母最小公倍数律]\label{cor:pom-anom-torus-phase-elimination}
\leanverified{paper\_pom\_anom\_torus\_phase\_elimination}
设异常取值于紧环面
\(
H=(\RR/2\pi\ZZ)^r
\)，并写 \(a=(a_1,\dots,a_r)\in H\)。
则存在整数 \(q\ge 2\) 使 \(qa=0\) 于 \(H\) 当且仅当对每个 \(1\le j\le r\) 都有
\[
a_j\in 2\pi\QQ/2\pi\ZZ.
\]
在该情形下，若对每个 \(j\) 取既约表示
\(
a_j\equiv 2\pi\,u_j/v_j\pmod{2\pi}
\)（\(\gcd(u_j,v_j)=1\)），则最小消去阶为
\[
\boxed{\;
q_\star=\mathrm{lcm}(v_1,\dots,v_r)\; },
\]
并且对任意整数 \(q\ge 1\) 有 \(qa=0\) 当且仅当 \(q_\star\mid q\)。
\end{corollary}
\begin{proof}
在 \(H\) 中有 \(qa=0\) 当且仅当对每个 \(j\) 都有 \(q a_j\equiv 0\ (\mathrm{mod}\ 2\pi)\)，亦即 \(a_j/(2\pi)\in \QQ/\ZZ\) 且其分母整除 \(q\)。
最小正整数由分母的最小公倍数给出。证毕。
\end{proof}

\begin{corollary}[单位根相位的 torsion 充分条件]\label{cor:pom-anom-cyclotomic-phase-torsion}
\leanverified{paper\_pom\_anom\_cyclotomic\_phase\_torsion}
在推论 \ref{cor:pom-anom-torus-phase-elimination} 的设定下，若存在正整数 \(N\) 使得对每个坐标 \(a_j\) 都满足 \(N a_j\equiv 0\ (\mathrm{mod}\ 2\pi)\)（等价于 \(e^{\mathrm{i}a_j}\) 为有限阶单位根），则 \(a\in H_{\mathrm{tor}}\)，从而存在有限阶 \(q\ge 2\) 使 \([\mathrm{Anom}(\mathrm{MOM}_q\circ w)]=0\)。
\end{corollary}
\begin{proof}
由假设可取同一 \(N\) 使 \(Na=0\) 于 \(H\)，故 \(a\) 为 torsion 元。
再由定理 \ref{thm:pom-anom-torsion-eliminability-min-order} 得结论。证毕。
\end{proof}

\begin{corollary}[非 torsion 异常的无界放大与阈值穿越]\label{cor:pom-anom-unbounded-amplification}
\leanverified{paper\_pom\_anom\_unbounded\_amplification}
在定理 \ref{thm:pom-anom-torsion-eliminability-min-order} 的设定下，若异常类 \(a=[\mathrm{Anom}(w)]\) 不是 torsion 元，则集合
\(
\{[\mathrm{Anom}(\mathrm{MOM}_q\circ w)]:q\ge 2\}
\)
在 \(H\) 中无穷。
进一步，若 \(H\) 可视为某个赋范实线性空间 \((V,\|\cdot\|)\) 的加法子群（例如 \(H\) 本身为实向量空间或其加法子群），并以该范数限制为 \(|\cdot|\)，则对任意阈值 \(\tau>0\) 存在整数 \(q\ge 2\) 使得
\[
\boxed{\;
\bigl|[\mathrm{Anom}(\mathrm{MOM}_q\circ w)]\bigr|>\tau\; }.
\]
\end{corollary}
\begin{proof}
由命题 \ref{prop:pom-anom-moment-amplification}，对 \(q\ge 2\) 有
\(
[\mathrm{Anom}(\mathrm{MOM}_q\circ w)]=qa
\)。
若存在 \(q_1\neq q_2\) 使 \(q_1a=q_2a\)，则 \((q_1-q_2)a=0\)，迫使 \(a\in H_{\mathrm{tor}}\)，与假设矛盾；故 \(\{qa:q\ge 2\}\) 无穷。
在赋范嵌入下，齐次性给出 \(|qa|=q|a|\to\infty\)，从而可取 \(q>\tau/|a|\) 得阈值穿越。证毕。
\end{proof}

\begin{proposition}[交叉异常的支撑约束：平凡坐标处必消失]\label{prop:pom-cross-anom-support-constraint}
\leanverified{paper\_pom\_cross\_anom\_support\_constraint}
在定义 \ref{def:pom-higher-cross-anomaly} 的设定下，对任意 \(|S|\ge 2\) 与任意 \(\chi_S\in\prod_{i\in S}\widehat{G_i}\)，若存在某个 \(i\in S\) 使 \(\chi_i=\chi_{i,0}\) 为平凡角色，则
\[
\boxed{\;
\Delta_S(\chi_S;\theta)=0\; }.
\]
\end{proposition}
\begin{proof}
若 \(\chi_i\) 在某坐标为平凡角色，则对定义 \ref{def:pom-higher-cross-anomaly} 的交替求和中任意 \(T\subseteq S\)，
项 \(\log\mathfrak{M}_{\chi^{(T)}}(\theta)\) 与 \(\log\mathfrak{M}_{\chi^{(T\cup\{i\})}}(\theta)\) 相同（因为第 \(i\) 坐标均取平凡角色），而它们在求和中系数符号相反，故成对抵消。
遍历所有包含/不包含 \(i\) 的配对即得 \(\Delta_S(\chi_S;\theta)=0\)。证毕。
\end{proof}

\begin{theorem}[Euler 缺陷的商空间坐标化：高阶交叉异常给出最小线性不变量]\label{thm:pom-euler-defect-quotient-coordinate}
固定 \(\theta\)，令 \(\widehat G=\prod_{i=1}^k\widehat{G_i}\)，并定义函数空间
\[
\mathcal{F}:=\{f:\widehat G\to\RR:\ f(\chi^0)=0\},
\qquad
\chi^0:=(\chi_{1,0},\dots,\chi_{k,0}).
\]
令“轴可分解”子空间
\[
\mathcal{S}:=
\Bigl\{f\in\mathcal{F}:\ \exists\,L_i:\widehat{G_i}\to\RR,\ L_i(\chi_{i,0})=0,\ f(\chi)=\sum_{i=1}^k L_i(\chi_i)\Bigr\}.
\]
对每个 \(|S|\ge 2\)，令 \(V_S\) 为 \(\prod_{i\in S}\widehat{G_i}\) 上的函数空间，并要求在任一坐标平凡时取零（命题 \ref{prop:pom-cross-anom-support-constraint}）。
定义线性映射
\[
\Phi:\mathcal{F}\longrightarrow \bigoplus_{|S|\ge 2}V_S,\qquad
\Phi(f):=\bigl(\Delta_S(\cdot;\theta)\bigr)_{|S|\ge 2},
\]
其中 \(\Delta_S\) 由定义 \ref{def:pom-higher-cross-anomaly}（对 \(f=\log\mathfrak{M}_\chi(\theta)\) 的一般化）给出。
则：
\begin{enumerate}
  \item \(\ker\Phi=\mathcal{S}\)。等价地，所有高阶交叉异常全消失当且仅当 \(f\) 轴可分解（定理 \ref{thm:pom-axis-decomposability-iff} 的线性化表述）。
  \item \(\Phi\) 诱导出商空间同构
  \[
  \boxed{\;
  \overline\Phi:\mathcal{F}/\mathcal{S}\xrightarrow{\ \sim\ }\bigoplus_{|S|\ge 2}V_S\; }.
  \]
  \item 每个商类 \([f]\in\mathcal{F}/\mathcal{S}\) 存在唯一代表 \(f^{\mathrm{int}}\) 满足全部一体项（单坐标差分）恒为零，即
  \(
  \Delta_{\{i\}}(\cdot;\theta)\equiv 0
  \)
  对一切 \(i\)，并且
  \[
  f^{\mathrm{int}}(\chi)=\sum_{|S|\ge 2}\Delta_S(\chi_S;\theta).
  \]
\end{enumerate}
\leanverified{paper\_pom\_euler\_defect\_quotient\_coordinate}
\end{theorem}
\begin{proof}
(1) 由定理 \ref{thm:pom-axis-decomposability-iff}，轴可分解当且仅当对一切 \(|S|\ge 2\) 有 \(\Delta_S\equiv 0\)，这正是 \(\ker\Phi=\mathcal{S}\)。
\par
(2) 由 (1) 知 \(\Phi\) 在商空间上诱导单射。
满射性可由显式构造给出：给定任意族 \((g_S)_{|S|\ge 2}\)（其中 \(g_S\in V_S\)），令
\(
f(\chi):=\sum_{|S|\ge 2} g_S(\chi_S)
\)。
由定义 \ref{def:pom-higher-cross-anomaly} 的差分表达（定理 \ref{thm:pom-mobius-complete-decomposition} 的证明思路），对任意 \(|S|\ge 2\) 计算 \(\Delta_S\) 时，只有包含于 \(S\) 的项可能贡献；而 \(g_T\in V_T\) 在任一坐标平凡时取零，使得 \(T\neq S\) 的项在 \(S\)-差分下抵消，最终得到 \(\Delta_S=g_S\)。
\par
(3) 对任意 \(f\in\mathcal{F}\)，由定理 \ref{thm:pom-mobius-complete-decomposition} 有
\(
f=\sum_{i}\Delta_{\{i\}}+\sum_{|S|\ge 2}\Delta_S
\)。
令
\(
f^{\mathrm{int}}:=f-\sum_i \Delta_{\{i\}}
\)，
则 \(f^{\mathrm{int}}\equiv f\ (\mathrm{mod}\ \mathcal{S})\) 且其一体项为零，并满足所示表示。
若另有代表 \(f'\) 亦满足 \(\Delta_{\{i\}}\equiv 0\)，则 \(f-f'\in\mathcal{S}\) 且其单坐标分量全为零，只能为零函数，故唯一性成立。证毕。
\end{proof}

\begin{corollary}[维数公式（有限阿贝尔轴）]\label{cor:pom-euler-defect-dimension-formula}
\leanverified{paper\_pom\_euler\_defect\_dimension\_formula}
在定理 \ref{thm:pom-euler-defect-quotient-coordinate} 的设定下，若各轴为有限阿贝尔群并记
\(
n_i:=|\widehat{G_i}|=|G_i|
\)，则
\[
\boxed{\;
\dim(\mathcal{F}/\mathcal{S})
\;=\;
\prod_{i=1}^k n_i - 1 - \sum_{i=1}^k (n_i-1)
\;=\;
\sum_{|S|\ge 2}\prod_{i\in S}(n_i-1)\; }.
\]
\end{corollary}
\begin{proof}
\(\dim\mathcal{F}=\prod_i n_i-1\)（归一化条件去掉一个自由度）。
而 \(\mathcal{S}\cong\bigoplus_{i=1}^k\{L_i:\widehat{G_i}\to\RR,\ L_i(\chi_{i,0})=0\}\)，故
\(\dim\mathcal{S}=\sum_i(n_i-1)\)。
最后一条恒等式来自
\(
\prod_i(1+(n_i-1))=\sum_{S\subseteq[k]}\prod_{i\in S}(n_i-1)
\)
并剔除 \(S=\varnothing\) 与单元素集合项。证毕。
\end{proof}

\begin{theorem}[高阶异常签名的 Möbius--ANOVA 分解及其精确维数]\label{thm:pom-cross-anom-mobius-anova-dimension-theory}
\leanverified{paper\_pom\_cross\_anom\_mobius\_anova\_dimension\_theory}
设 \(G_1,\dots,G_k\) 为有限阿贝尔群，记
\[
\widehat G:=\widehat G_1\times\cdots\times\widehat G_k,
\qquad
\chi^0:=(\chi_{1,0},\dots,\chi_{k,0}),
\]
并定义归一化函数空间
\[
\mathcal A_k:=\{f:\widehat G\to\RR:\ f(\chi^0)=0\}.
\]
对任意非空子集 \(S\subseteq[k]\)，定义 \(\mathcal A_S\subset \mathcal A_k\) 为所有满足下列条件的函数：
\begin{enumerate}
  \item \(f(\chi)\) 仅依赖于子元组 \(\chi_S=(\chi_i)_{i\in S}\)；
  \item 若某个 \(i\in S\) 满足 \(\chi_i=\chi_{i,0}\)，则 \(f(\chi)=0\)。
\end{enumerate}
则存在规范直和分解
\[
\boxed{
\mathcal A_k=\bigoplus_{\varnothing\neq S\subseteq[k]}\mathcal A_S.
}
\]
并且每个纯 \(S\)-体分量都满足精确维数公式
\[
\boxed{
\dim \mathcal A_S=\prod_{i\in S}\bigl(|G_i|-1\bigr).
}
\]
因而总维数为
\[
\boxed{
\dim\mathcal A_k
=
\sum_{\varnothing\neq S\subseteq[k]}\prod_{i\in S}\bigl(|G_i|-1\bigr)
=
\prod_{i=1}^{k}|G_i|-1.
}
\]
特别地，只含一体项的轴可分解子空间维数为
\[
\sum_{i=1}^{k}(|G_i|-1),
\]
而所有两体及以下可见分量的总维数为
\[
\sum_{i=1}^{k}(|G_i|-1)
+
\sum_{1\le i<j\le k}(|G_i|-1)(|G_j|-1).
\]
因此三体及更高阶隐藏耦合的精确余维数为
\[
\sum_{|S|\ge 3}\prod_{i\in S}(|G_i|-1).
\]
\end{theorem}
\begin{proof}
对任意 \(f\in\mathcal A_k\) 与非空 \(S\subseteq[k]\)，定义 Möbius 分量
\[
f_S(\chi)
:=
\sum_{T\subseteq S}(-1)^{|S|-|T|}f(\chi^{(T)}),
\]
其中 \(\chi^{(T)}\) 表示在 \(T\) 上保留原坐标、其余坐标取平凡角色的元组。由定理 \ref{thm:pom-mobius-complete-decomposition} 的同一 Möbius 反演计算，
\[
f=\sum_{\varnothing\neq S\subseteq[k]} f_S.
\]
并且每个 \(f_S\) 仅依赖于 \(\chi_S\)，且若某个 \(i\in S\) 上 \(\chi_i=\chi_{i,0}\)，则交替和按包含/不包含 \(i\) 的项成对抵消，故 \(f_S\in \mathcal A_S\)。由 Möbius 反演的唯一性，上述分解是规范且唯一的，从而得到直和分解。

现算维数。固定非空 \(S\subseteq[k]\)。由于 \(\mathcal A_S\) 中的函数在任一平凡坐标处必为零，故其全部自由参数恰落在
\[
\prod_{i\in S}\bigl(\widehat G_i\setminus\{\chi_{i,0}\}\bigr)
\]
上；反之，在该集合上任意指定实值后，都可唯一扩张为 \(\mathcal A_S\) 中的函数。于是
\[
\dim \mathcal A_S=\prod_{i\in S}\bigl(|\widehat G_i|-1\bigr)
=
\prod_{i\in S}\bigl(|G_i|-1\bigr).
\]
对所有非空 \(S\) 求和即得
\[
\dim\mathcal A_k
=
\sum_{\varnothing\neq S\subseteq[k]}\prod_{i\in S}(|G_i|-1).
\]
再由恒等式
\[
\prod_{i=1}^{k}\bigl(1+(|G_i|-1)\bigr)
=
\sum_{S\subseteq[k]}\prod_{i\in S}(|G_i|-1)
\]
并剔除空集项，即得
\[
\dim\mathcal A_k=\prod_{i=1}^{k}|G_i|-1.
\]
最后，把 \(|S|=1\)、\(|S|\le 2\)、\(|S|\ge 3\) 的项分别分组，便得到轴可分解维数、两体可见维数与高阶隐藏余维数公式。证毕。
\end{proof}

\begin{corollary}[异常签名塔与对称幂 moment-kernel 的指数--多项式复杂度分离]\label{cor:pom-anom-signature-tower-sympower-complexity-separation}
\leanverified{paper\_pom\_anom\_signature\_tower\_sympower\_complexity\_separation}
设 \(G\) 为任意非平凡有限阿贝尔群，\(|G|\ge 2\)。对每个整数 \(q\ge 2\)，定义 \ref{def:pom-anom-signature-tower} 的 \(q\)-阶异常签名塔
\[
\mathsf{Anom}^{(q)}_G
\]
的坐标数精确等于
\[
\boxed{
|G|^{q-1}-1.
}
\]
另一方面，若底层在线同步核的状态数为 \(|Q|\)，则由命题 \ref{prop:pom-ak-symmetric-compression}，\(q\)-重碰撞对称幂商核的状态数至多为
\[
\boxed{
\binom{q+|Q|-1}{|Q|-1}=O\!\left(q^{|Q|-1}\right).
}
\]
因而
\[
\boxed{
\frac{|G|^{q-1}-1}{\binom{q+|Q|-1}{|Q|-1}}
\xrightarrow[q\to\infty]{}
\infty.
}
\]
这说明异常语义侧的证书维数按 \(q\) 指数增长，而对称幂 moment-kernel 的最优状态压缩至多保持多项式增长；二者之间存在不可消除的指数--多项式分离。
\end{corollary}
\begin{proof}
由定义 \ref{def:pom-anom-signature-tower}，满足约束
\[
\chi_1\cdots\chi_q=\chi_0
\]
的 \(q\)-元角色组中，前 \(q-1\) 个角色可以在 \(\widehat G\) 中任意选取，而最后一个由乘积约束唯一决定，因此总数恰为 \(|\widehat G|^{q-1}=|G|^{q-1}\)。去掉全平凡元组 \((\chi_0,\dots,\chi_0)\) 后，即得坐标数
\[
|G|^{q-1}-1.
\]
另一方面，命题 \ref{prop:pom-ak-symmetric-compression} 已给出 \(q\)-重 moment-kernel 的对称幂商状态数上界
\[
\dim(\widetilde A_q)\le \binom{q+|Q|-1}{|Q|-1}.
\]
当 \(|Q|\) 固定时，上式是 \(q\) 的 \((|Q|-1)\) 次多项式，而 \(|G|^{q-1}-1\) 以底数 \(|G|>1\) 指数增长，故其比值趋于 \(+\infty\)。证毕。
\end{proof}

\begin{theorem}[Euler 缺陷的正交分解：交叉项能量的 Pythagoras 恒等式]\label{thm:pom-euler-defect-orthogonal-pythagoras}
\leanverified{paper\_pom\_euler\_defect\_orthogonal\_pythagoras}
在定理 \ref{thm:pom-euler-defect-quotient-coordinate} 的设定下，赋予 \(\widehat G=\prod_{i=1}^k\widehat{G_i}\) 以均匀测度，并将 \(\mathcal F\subset L^2(\widehat G)\) 视为 Hilbert 子空间。
对任意 \(f\in\mathcal F\)，定义坐标平均算子
\[
(E_i f)(\chi):=\mathbb{E}_{\xi_i\sim \mathrm{Unif}(\widehat{G_i})} f(\chi_1,\dots,\chi_{i-1},\xi_i,\chi_{i+1},\dots,\chi_k),
\qquad
D_i:=I-E_i.
\]
对任意子集 \(S\subseteq[k]\) 定义
\(
D_S:=\prod_{i\in S}D_i
\)、
\(
E_{S^c}:=\prod_{i\notin S}E_i
\)，并置
\[
f_S:=E_{S^c}D_S f.
\]
则：
\begin{enumerate}
  \item 存在严格正交分解
  \[
  \boxed{\;
  f=\sum_{S\subseteq[k]} f_S,\qquad
  \langle f_S,f_T\rangle_{L^2}=0\ (S\ne T)\;}
  \]
  并且满足 Pythagoras 恒等式
  \[
  \|f\|_{L^2}^2=\sum_{S\subseteq[k]}\|f_S\|_{L^2}^2.
  \]
  \item 对每个 \(\emptyset\ne S\subseteq[k]\) 与每个 \(i\in S\)，有零均值条件 \(E_i f_S=0\)；并且 \(f_{\varnothing}=E_{[k]}f\) 为常数函数（等于 \(f\) 的全局均值）。
\end{enumerate}
\end{theorem}
\begin{proof}
由于 \(E_i\) 为 \(L^2\) 上对“与第 \(i\) 轴无关”的子空间的条件期望投影，故 \(E_i\) 两两交换且自伴幂等；从而 \(D_i=I-E_i\) 亦两两交换且自伴。
展开恒等式
\[
I=\prod_{i=1}^k(E_i+D_i)=\sum_{S\subseteq[k]}E_{S^c}D_S
\]
并作用于 \(f\) 得到分解式。
若 \(S\ne T\)，取某个 \(i\in S\triangle T\)，则可将 \(f_S\) 写为 \(D_i(\cdot)\) 的像、而 \(f_T\) 写为 \(E_i(\cdot)\) 的像；由 \(E_i\) 为正交投影（等价于 \(\langle D_i g, E_i h\rangle_{L^2}=0\)）得 \(\langle f_S,f_T\rangle_{L^2}=0\)，从而得到正交与 Pythagoras 恒等式。
而对 \(i\in S\)，由定义 \(f_S\) 含有因子 \(D_i\)，故 \(E_i f_S=E_iD_i(\cdot)=0\)。证毕。
\end{proof}

\begin{theorem}[最小作用量 gauge：交叉项能量等于到可加子空间的最小距离]\label{thm:pom-euler-defect-min-action-gauge}
\leanverified{paper\_pom\_euler\_defect\_min\_action\_gauge}
沿用定理 \ref{thm:pom-euler-defect-orthogonal-pythagoras} 的记号，并定义可加（轴可分解）子空间
\[
\mathcal S_{\mathrm{add}}
:=
\Bigl\{g\in L^2(\widehat G):\ \exists\,c\in\RR,\ \exists\,g_i:\widehat{G_i}\to\RR,\ g_i(\chi_{i,0})=0,\ g(\chi)=c+\sum_{i=1}^k g_i(\chi_i)\Bigr\}.
\]
则 \(L^2\) 意义下到 \(\mathcal S_{\mathrm{add}}\) 的正交投影为
\[
\boxed{\ \Pi_{\mathcal S_{\mathrm{add}}} f=f_{\varnothing}+\sum_{|S|=1} f_S\ }.
\]
并且最小残差能量满足
\[
\boxed{\ \inf_{g\in\mathcal S_{\mathrm{add}}}\|f-g\|_{L^2}^2=\sum_{|S|\ge 2}\|f_S\|_{L^2}^2\ }.
\]
右端因此给出一条规范化的“Euler 缺陷能量”：其为零当且仅当 \(f\in\mathcal S_{\mathrm{add}}\)（无交叉项）。
此外，定理 \ref{thm:pom-euler-defect-quotient-coordinate} 的轴可分解子空间满足 \(\mathcal S=\mathcal S_{\mathrm{add}}\cap\mathcal F\)。
\end{theorem}
\begin{proof}
由定理 \ref{thm:pom-euler-defect-orthogonal-pythagoras}，\(f\) 的 Hoeffding 分量 \(\{f_S\}_{S\subseteq[k]}\) 两两正交。
记
\[
\mathcal H_{\le 1}
:=\mathrm{Im}(E_{[k]})\ \oplus\ \bigoplus_{i=1}^k \mathrm{Im}(E_{[k]\setminus\{i\}}D_i),
\]
其中 \(\mathrm{Im}(E_{[k]})\) 为常数函数子空间。
由定义可知 \(\mathcal H_{\le 1}\subseteq \mathcal S_{\mathrm{add}}\)，并且对任意 \(g\in\mathcal S_{\mathrm{add}}\)，其二阶及以上交叉分量必为零（即对所有 \(|S|\ge 2\) 有 \(g_S=0\)），从而 \(g=\sum_{|S|\le 1} g_S\in \mathcal H_{\le 1}\)。
因此 \(\mathcal S_{\mathrm{add}}=\mathcal H_{\le 1}\)。
故 \(f_{\varnothing}+\sum_{|S|=1}f_S\) 即为 \(f\) 在闭子空间 \(\mathcal S_{\mathrm{add}}\) 上的正交投影，而最小距离平方由正交补能量之和给出，即
\(
\sum_{|S|\ge 2}\|f_S\|_{L^2}^2
\)。
最后，“能量为零”当且仅当所有 \(|S|\ge 2\) 的分量消失，亦即 \(f\in\mathcal S_{\mathrm{add}}\)。证毕。
\end{proof}
